\section{Methods}

Full methods have been published \cite{laws2025modelling}.

\subsection{Lower Super Output Areas (LSOA)}

We performed our analysis across England using LSOA as our geographic footprint \cite{ons}. LSOAs are small areas in England, with roughly 1,500 people in each area. There are 32,843 LSOAs in England.

\subsection{Travel times}

Travel times between all Lower Super Output Areas (LSOAs, areas with a population of approximately 1,500) and all stroke units were estimated using Open Street Map data \footnote{\url{https://www.openstreetmap.org/}} and the Routino \footnote{\url{https://www.routino.org/}} route finding engine on Linux.

\subsection{Outcome modelling}

Outcome modelling is based on published meta-analysis of clinical trials which have evaluated the decline of effectiveness  of IVT and MT over time \cite{emberson_effect_2014, fransen_time_2016}. We have applied those models to the expected real world treatment population (based on the national stroke audit of England and Wales, SSNAP). We calculated the expected benefit in the treated population for non large vessel occlusion (\textit{nLVO}, receiving IVT only) and large vessel occlusion (\textit{LVO}, receiving IVT and/or MT) separately.

We used modified Rankin Scale (mRS) at 3--6 months as a measure of outcome. mRS is the most commonly used instrument to describe post-stroke functional outcome \cite{quinn_functional_2009}, describing independence of living from a scale of 0 (no disability) through to 5 (severe disability requiring constant nursing attention), with death assigned an mRS of 6. A commonly used surrogate for independent living is mRS 0--2.

We calculated the patients' mRS outcome distribution based on time to treatment for three patient-treatment cohorts: nLVO treated with IVT; LVO treated with IVT alone; and LVO treated with IVT and MT. For each patient-treatment cohort we calculated an mRS distribution for treatment at any given time by interpolating between the mRS distribution for treatment given at \emph{t=0} (time of stroke onset) and the mRS distribution for treatment given at \emph{t=No Effect} (time of no effect of treatment), assuming that log odds of achieving any given mRS threshold fall linearly over time \cite{emberson_effect_2014, fransen_time_2016}. As possible changes to service provision reported here should only affect potential candidates for thrombectomy, in this report we only show results for treated large vessel occlusion stroke.

The time to no effect was 6.3 hours for IVT \cite{emberson_effect_2014} and 8.0 hours for MT \cite{ fransen_time_2016}. Our model did not include selection of patients who may still benefit from treatment beyond these durations through the use of perfusion scanning. This number is small for IVT, but is more substantial for MT – approximately 2500 per annum in England.

Our modelling compares outcomes for those patients receiving IVT and MT. We do not make any assumption about the proportion of all patients receiving those treatments.

Code for outcome modelling, and a further description of methods is available on GitHub \footnote{\url{https://github.com/stroke-modelling/stroke-outcome/}}, and we have made it available as a Python package on PyPi \footnote{\url{https://pypi.org/project/stroke-outcome/}}. We have described the methods in detail in an online book \footnote{\url{https://samuel-book.github.io/stroke_outcome/intro.html}}.

\subsection{Process times}

To study geographic variation in benefit of MSU care, a single set of parameters was chosen that reflected a reasonable base case for performance of usual care of LVO stroke. Process times (minutes) are shown in the following list.

\vspace{5mm}

\begin{minipage}{1.0\textwidth}  % Define the width of the minipage
\begin{itemize}
    \item Stroke onset to call: 30
    \item Call to ambulance arrival: 20
    \item Ambulance on-scene: 30
    \item Travel time to first admitting unit: Depends on geography
    \item Hospital arrival to IVT: 30
    \item If transfer for MT required:
    \begin{itemize}
        \item Transfer-related net MT delay (excluding travel time)*: 60
        \item Transport time to MT unit: Depends on geography
    \end{itemize}
    \item Hospital arrival to MT: 90
\end{itemize}

*This assumes a typical door-in-door-out time of 90 mints at the IVT-only centre, but 1 30 minute reduction in Hospital arrival to MT at the transferred-in MT centre, due to work already performed at the IVT-only unit.
\end{minipage}


